% GENERATED FILE. Edit canonical lessons, then run npm run generate:books.
% canonical-source-hash: fnv1a64:3579c41cc33f7b2c
% canonical-lessons: GU-C175-punam, GU-C175-amas, GU-C175-ant, GU-C175-addhu, GU-C175-satat

\chapter{Full moon and new moon}
\label{ch:gu-full-moon-and-new-moon}
\hlchaptermodality{175}

\begin{chapteropening}
\textbf{By the end of this chapter:} \emph{I can talk about the full moon, the new moon and the end of the month.}

Complete the last lesson of chapter 175: I can talk about the full moon, the new moon and the end of the month.
\end{chapteropening}

\section[\texorpdfstring{\gu{પૂનમ} (pūnam)}{પૂનમ (pūnam)}]{\gu{પૂનમ} (pūnam) — full moon day}

\label{lesson:GU-C175-punam}



\begin{quote}
Before the new one: say the Gujarati for wall calendar, then the Gujarati for beginning.
\end{quote}

\subsection*{You'll want to know: \gu{પૂનમ}}
\textbf{\gu{પૂનમ}} — \emph{pūnam} — \textquotedblleft{}full moon day\textquotedblright{}.

\begin{quote}
\textbf{\gu{પૂનમે} \gu{અમે} \gu{મંદિરે} \gu{જઈશું}.} — \emph{pūname ame mandire jaīśũ} — \textquotedblleft{}We will go to the temple on the full moon day.\textquotedblright{}
\end{quote}

\begin{cousinweb}
\textbf{\gu{પૂનમ}} (\emph{pūnam}) is usually traced to Sanskrit \emph{pūrṇimā}, \textquotedblleft{}full-moon day\textquotedblright{}, and EtymDB derives it from that word, but the path is not settled: CDIAL lists the Sanskrit word with no Gujarati form under it (CDIAL 8340).
\end{cousinweb}

\begin{grammarlens}[title={What you've built}]
One more word for everyday life at A2.
\end{grammarlens}

\subsection*{Your turn}
\begin{itemize}
\raggedright
  \item \emph{Say it:} \emph{pūnam}
  \item \emph{Say it:} \emph{pūnam}, once more
  \item \emph{Say it:} \emph{pūname ame mandire jaīśũ} — We will go to the temple on the full moon day.
  \item \emph{Recall:} say the Gujarati for wall calendar, then the Gujarati for beginning, then say \emph{pūnam} again
\end{itemize}

\begin{tcolorbox}[breakable,colback=teal!4,colframe=teal!35!black,title={Before you move on}]
What does \gu{પૂનમ} mean? (\textquotedblleft{}Full moon day\textquotedblright{}.) Say it once more.
\end{tcolorbox}
\section[\texorpdfstring{\gu{અમાસ} (amās)}{અમાસ (amās)}]{\gu{અમાસ} (amās) — new moon day}

\label{lesson:GU-C175-amas}



\begin{quote}
Before the new one: say the Gujarati for beginning, then the Gujarati for full moon day.
\end{quote}

\subsection*{You'll want to know: \gu{અમાસ}}
\textbf{\gu{અમાસ}} — \emph{amās} — \textquotedblleft{}new moon day\textquotedblright{}.

\begin{quote}
\textbf{\gu{અમાસની} \gu{રાત્રે} \gu{અંધારું} \gu{હશે}.} — \emph{amāsnī rātre andhārũ haśe} — \textquotedblleft{}It will be dark on the night of the new moon.\textquotedblright{}
\end{quote}

\begin{cousinweb}
\textbf{\gu{અમાસ}} (\emph{amās}) is a tadbhava. It comes from Sanskrit \emph{amāvāsyā}, \textquotedblleft{}the night of the new moon\textquotedblright{}, and CDIAL lists the Gujarati \emph{amās} with that sense (CDIAL 565).
\end{cousinweb}

\begin{grammarlens}[title={What you've built}]
One more word for everyday life at A2.
\end{grammarlens}

\subsection*{Your turn}
\begin{itemize}
\raggedright
  \item \emph{Say it:} \emph{amās}
  \item \emph{Say it:} \emph{amās}, once more
  \item \emph{Say it:} \emph{amāsnī rātre andhārũ haśe} — It will be dark on the night of the new moon.
  \item \emph{Recall:} say the Gujarati for beginning, then the Gujarati for full moon day, then say \emph{amās} again
\end{itemize}

\begin{tcolorbox}[breakable,colback=teal!4,colframe=teal!35!black,title={Before you move on}]
What does \gu{અમાસ} mean? (\textquotedblleft{}New moon day\textquotedblright{}.) Say it once more.
\end{tcolorbox}
\section[\texorpdfstring{\gu{અંત} (ant)}{અંત (ant)}]{\gu{અંત} (ant) — end}

\label{lesson:GU-C175-ant}



\begin{quote}
Before the new one: say the Gujarati for full moon day, then the Gujarati for new moon day.
\end{quote}

\subsection*{You'll want to know: \gu{અંત}}
\textbf{\gu{અંત}} — \emph{ant} — \textquotedblleft{}end\textquotedblright{}.

\begin{quote}
\textbf{\gu{મહિનાના} \gu{અંતમાં} \gu{હું} \gu{આવીશ}.} — \emph{mahinānā antmā̃ hũ āvīś} — \textquotedblleft{}I will come at the end of the month.\textquotedblright{}
\end{quote}

\begin{cousinweb}
\textbf{\gu{અંત}} (\emph{ant}) is Sanskrit \emph{anta}, \textquotedblleft{}end, border\textquotedblright{}, taken over whole as a learned loan, a tatsama. CDIAL lists the Sanskrit word (CDIAL 347).
\end{cousinweb}

\begin{grammarlens}[title={What you've built}]
One more word for everyday life at A2.
\end{grammarlens}

\subsection*{Your turn}
\begin{itemize}
\raggedright
  \item \emph{Say it:} \emph{ant}
  \item \emph{Say it:} \emph{ant}, once more
  \item \emph{Say it:} \emph{mahinānā antmā̃ hũ āvīś} — I will come at the end of the month.
  \item \emph{Recall:} say the Gujarati for full moon day, then the Gujarati for new moon day, then say \emph{ant} again
\end{itemize}

\begin{tcolorbox}[breakable,colback=teal!4,colframe=teal!35!black,title={Before you move on}]
What does \gu{અંત} mean? (\textquotedblleft{}End\textquotedblright{}.) Say it once more.
\end{tcolorbox}
\section[\texorpdfstring{\gu{અડધું} (aḍdhũ)}{અડધું (aḍdhũ)}]{\gu{અડધું} (aḍdhũ) — half}

\label{lesson:GU-C175-addhu}



\begin{quote}
Before the new one: say the Gujarati for new moon day, then the Gujarati for end.
\end{quote}

\subsection*{You'll want to know: \gu{અડધું}}
\textbf{\gu{અડધું}} — \emph{aḍdhũ} — \textquotedblleft{}half\textquotedblright{}.

\begin{quote}
\textbf{\gu{હું} \gu{અડધી} \gu{રોટલી} \gu{ખાઈશ}.} — \emph{hũ aḍdhī roṭlī khāīś} — \textquotedblleft{}I will eat half a rotli.\textquotedblright{}
\end{quote}

\begin{cousinweb}
\textbf{\gu{અડધું}} (\emph{aḍdhũ}) is inherited from Sanskrit \emph{ardha}, \textquotedblleft{}half\textquotedblright{}. CDIAL lists Gujarati \emph{adh}, \textquotedblleft{}half\textquotedblright{}, in compounds under that word, but not \emph{aḍdhũ} itself, so the exact path is not settled (CDIAL 644).
\end{cousinweb}

\begin{grammarlens}[title={What you've built}]
One more word for everyday life at A2.
\end{grammarlens}

\subsection*{Your turn}
\begin{itemize}
\raggedright
  \item \emph{Say it:} \emph{aḍdhũ}
  \item \emph{Say it:} \emph{aḍdhũ}, once more
  \item \emph{Say it:} \emph{hũ aḍdhī roṭlī khāīś} — I will eat half a rotli.
  \item \emph{Recall:} say the Gujarati for new moon day, then the Gujarati for end, then say \emph{aḍdhũ} again
\end{itemize}

\begin{tcolorbox}[breakable,colback=teal!4,colframe=teal!35!black,title={Before you move on}]
What does \gu{અડધું} mean? (\textquotedblleft{}Half\textquotedblright{}.) Say it once more.
\end{tcolorbox}
\section[\texorpdfstring{\gu{સતત} (satat)}{સતત (satat)}]{\gu{સતત} (satat) — constantly, nonstop}

\label{lesson:GU-C175-satat}



\begin{quote}
Before the new one: say the Gujarati for end, then the Gujarati for half.
\end{quote}

\subsection*{You'll want to know: \gu{સતત}}
\textbf{\gu{સતત}} — \emph{satat} — \textquotedblleft{}constantly, nonstop\textquotedblright{}.

\begin{quote}
\textbf{\gu{તે} \gu{સતત} \gu{કામ} \gu{કરશે}.} — \emph{te satat kām karśe} — \textquotedblleft{}He will work nonstop.\textquotedblright{}
\end{quote}

\begin{cousinweb}
\textbf{\gu{સતત}} (\emph{satat}) is Sanskrit \emph{satata}, \textquotedblleft{}constant, uninterrupted\textquotedblright{}, taken over whole as a learned loan, a tatsama; EtymDB lists it as borrowed from Sanskrit. CDIAL lists the Sanskrit word but no Gujarati form (CDIAL 13104).
\end{cousinweb}

\begin{grammarlens}[title={What you've built}]
One more word for everyday life at A2.
\end{grammarlens}

\subsection*{Your turn}
\begin{itemize}
\raggedright
  \item \emph{Say it:} \emph{satat}
  \item \emph{Say it:} \emph{satat}, once more
  \item \emph{Say it:} \emph{te satat kām karśe} — He will work nonstop.
  \item \emph{Recall:} say the Gujarati for end, then the Gujarati for half, then say \emph{satat} again
\end{itemize}

\begin{tcolorbox}[breakable,colback=teal!4,colframe=teal!35!black,title={Before you move on}]
What does \gu{સતત} mean? (\textquotedblleft{}Constantly, nonstop\textquotedblright{}.) Say it once more.
\end{tcolorbox}
